\documentclass[a4paper]{article}
\usepackage[pdftex]{graphicx}
\usepackage[utf8]{inputenc}
\usepackage{enumerate}
\usepackage{siunitx}
\sisetup{locale=DE}
\usepackage{href-ul}
\hypersetup{
	colorlinks=true,
	linkcolor=blue,
	urlcolor=blue}
\usepackage{icomma}
\usepackage{amssymb}
\usepackage{geometry}
\geometry{a4paper, top=15mm, left=15mm, right=15mm, bottom=15mm,
	headsep=10mm, footskip=12mm}

\begin{document}
	%Title
	{\bf Capacitor in AC circuit }\\
	\begin{enumerate}[1.]
		
		%Task 1
		\begin{minipage}{0.65\textwidth}
			\item The input voltage is \
			$U(t)=\hat{U}\cdot \sin{(2 \pi f t)}$. The effective value $U_{eff}$ of the input voltage $U$ is constant and is \\ $U_{eff}=\SI{12}{\volt}$. The frequency of this input voltage $U$ can be changed.
			The output voltage can be the voltage $U_R$ at the ohmic resistor $R$ or the voltage at the capacitor. The effective values $U_{R,eff}$ and $U_{C,eff}$ depend on the frequency $f$ of the input voltage $U$.
		\end{minipage}
		\hfill
		\begin{minipage}{0.35\textwidth}
			\includegraphics[width=5.5 cm]{weko234}
		\end{minipage}
		
		Using a voltage phasor diagram, show that the impedance $Z$ is: $$ Z= \sqrt{R^2 + {1 \over 4 \pi^2 C^2} \cdot {1 \over f^ 2}}$$
		
		%Exercise 2
		\item The frequency $f$ of the input voltage $U$ is initially set to the value $f_1=\SI{520}{\hertz}$. At this frequency $f_1$, a current with the strength $I_1$ flows in the alternating current circuit and the voltage $U_{R,1}$ dropped across the ohmic resistor $R$ has the effective value
		$U_{R,eff,1}=\SI{11,2}{\volt}$
		\begin{enumerate}[a)]
			\item Calculate the capacitance $C$ of the capacitor and the phase shift $\varphi_1$ occurring between the current $I_1$ and the input voltage. $[\textnormal{Partial result:}\quad C=\SI{20}{\nano\farad}]$
			\item Calculate the average active power $P_1$ of the AC circuit.
		\end{enumerate}
		
		%Task 3
		\item The frequency $f$ of the input voltage $U$ is now varied. First, the voltage $U_R$ is tapped at the ohmic resistor as the output voltage.
		\begin{enumerate}[a)]
			\item Confirm by general calculation that for the effective value $U_{R,eff}$ of the voltage $U_R$, which depends on the frequency $f$, the following applies:
			$$U_{R,eff}(f)={R \over \sqrt{R^2 + {1 \over 4 \pi^2 C^2} \cdot {1 \over f^2}}}\cdot U_{eff}$$
			\item If the frequency $f$ of the input voltage $U$ is set to the value $f_g$, the effective values $U_{R,eff}$ and $U_{C,eff}$ are the same size. Calculate the so-called cutoff frequency $f_g$ and the effective values $U_{R,eff}$ and $U_{C,eff}$ for this frequency $f_g$. $[\textnormal{Partial result:}\quad f_g=\SI{0.20}{\kilo\hertz}]$
			\item Using the formula from subtask 3. a), justify that for high frequencies
			$f (f \gg f_g)$ the effective value $U_{R,eff}$ of the output voltage $U_R$ is practically as large as the effective value $U_{eff}$ of the input voltage $U$.
		\end{enumerate}
	\end{enumerate}
	
	\fbox{
		\begin{minipage}{0.5\textwidth}
			For the solution, please \href{https://www.okuyakl.de/phys/p13wekoL234/le234.pdf}{click here} or scan the QR code.\\
			Additional worksheets are available at
			
			\href{http://www.okuyakl.com}{www.okuyakl.com}
		\end{minipage}
		\hfill
		\begin{minipage}{0.4\textwidth}
			\includegraphics[width=1.5 cm]{../../viecher/zwe03}
			\includegraphics[width=3 cm]{qre234}
			\includegraphics[width=2 cm]{../../viecher/afanticon1}
			
	\end{minipage}}
	
\end{document}